\section{Memory Wall: el modelo crece más rápido que la capacidad de suministrar datos}

La sección anterior estableció las cotas inferiores de recursos; esta sección se centra en la memory wall, que se enfatizó repetidamente en clase. Gholami et al., en \emph{AI and Memory Wall}, analizan el desajuste entre el crecimiento del tamaño de los modelos, la memory capacity, el bandwidth y la comunicación. En un LLM, el training requiere parámetros, optimizer state, gradient y activation; la inference, además, debe gestionar la KV cache. Que el modelo «quepa» y que el modelo «reciba datos suficientes» son dos condiciones distintas.

\subsection{Muro de capacidad, muro de ancho de banda y muro de comunicación}

Esta sección descompone la «memory wall» en tres problemas medibles: si el modelo y su estado de ejecución caben, si pueden entregarse a compute por unidad de tiempo y si, al escalar a varias tarjetas, pueden sincronizarse de forma eficiente. Al leer la figura siguiente, no hay que agrupar el lado derecho bajo «falta memory», sino registrar por separado la evidencia de capacity, bandwidth y communication, porque cada tipo de problema requiere soluciones de hardware y software distintas.

\lecturefigure{03-memory-wall.png}{La esencia de la memory wall es el desajuste de velocidades entre compute throughput y data supply.}{Subtítulos locales 08:20--10:30; Gholami et al., \emph{AI and Memory Wall}.}

\begin{knowledgebox}{Lectura de la figura: primero, distinguir los tres tipos de wall}
La \textbf{capacity wall} pregunta si los parámetros, la KV cache y las activations caben en memory; la \textbf{bandwidth wall} pregunta si los datos pueden entregarse a compute por unidad de tiempo; la \textbf{communication wall} pregunta si la sincronización y la redistribución entre varios dispositivos absorben las ganancias del escalamiento. Las tres pueden presentarse al mismo tiempo, y resolver una de ellas expone el cuello de botella de la siguiente capa.
\end{knowledgebox}

\begin{importantbox}{La arithmetic intensity determina si el límite es «cómputo o ancho de banda»}
Se define la arithmetic intensity
\[
I=\frac{F}{B_m},
\]
es decir, cuántas operaciones se realizan por cada byte transferido. Si $I$ es muy baja, aunque el throughput pico del accelerator sea muy alto, es más probable que el kernel esté limitado por el memory bandwidth; el valor de fusion, tiling, quantization y cache reuse puede entenderse como aumentar la $I$ efectiva o reducir $B_m$.
\end{importantbox}

\subsection{Términos clave: SRAM, DRAM, HBM y storage}

Los subtítulos y los materiales comerciales suelen agrupar estos términos bajo «memory». El lector debe comprender primero el papel de cada uno en la hierarchy y, después, las diferencias entre Groq, NPU, GPU o las arquitecturas nuevas.

\begin{knowledgebox}{Concepto previo: cuatro niveles de almacenamiento}
\begin{tabularx}{\textwidth}{L{0.17\textwidth}L{0.24\textwidth}X}
\toprule
Término & Características principales & Papel y limitaciones en sistemas de IA \\
\midrule
SRAM & Static Random-Access Memory, caché de alta velocidad en el chip; baja latencia, pero costosa en área & Adecuada para datos pequeños y de uso frecuente, y para accesos deterministas; su capacidad es reducida y no equivale a un almacenamiento general para modelos. \\
DRAM & Dynamic Random-Access Memory; requiere refresco y tiene mayor capacidad & Memoria principal o memory externa del accelerator; su latencia y su ancho de banda dependen de la interfaz, el encapsulado y los accesos concurrentes. \\
HBM & High Bandwidth Memory, DRAM apilada de alto ancho de banda & Ofrece alto throughput mediante una interfaz ancha; la capacidad, el suministro de encapsulado, el costo y el consumo de energía siguen siendo restricciones. \\
SSD & Almacenamiento persistente, de gran capacidad pero con latencia mucho mayor que la memory & Guarda checkpoints, artifacts de modelos y datos fríos; normalmente requiere staging y no puede sustituir directamente a la HBM. \\
\bottomrule
\end{tabularx}
\end{knowledgebox}

\lecturefigure{04-memory-hierarchy.png}{Los distintos memory tiers intercambian latency, bandwidth, capacity y energy.}{Subtítulos locales 09:00--11:10; redibujo conceptual.}

\begin{knowledgebox}{Lectura de la figura: en el serving de modelos, los datos se ubican según su «temperatura»}
El kernel state más usado y los buffers pequeños permanecen en la SRAM on-chip; los weights activos y la KV cache compiten por la HBM; la host memory, más grande, se encarga del staging; el SSD guarda los modelos fríos y los checkpoints. Si el routing cambia de modelo con frecuencia, el arranque en frío y la transferencia de datos se convierten en el costo principal; por ello, el diseño del model portfolio debe coordinarse con la memory placement.
\end{knowledgebox}

\begin{warningbox}{«Tal chip solo sirve para texto» no es una conclusión estable}
En clase se usaron empresas concretas para ilustrar rápidamente las compensaciones entre SRAM, DRAM y HBM, pero las capacidades de los productos evolucionan con el compiler, el model support, la memory topology y el serving software. Estas notas conservan los principios de la hierarchy y no toman los juicios expresados en clase en 2025 como una clasificación permanente de productos.
\end{warningbox}

\subsection{Resumen de la sección}

La memory wall no es la escasez de un solo componente, sino un problema combinado de capacity, bandwidth, communication y software locality. Antes de elegir un chip, conviene cuantificar el working set del modelo, los bytes per token, el crecimiento de la KV cache y el traffic entre dispositivos.
